\documentclass[11pt]{article}

\usepackage[margin=1in]{geometry}
\usepackage{amsmath,amssymb}
\usepackage{bm}
\usepackage[hidelinks]{hyperref}

\newcommand{\kv}{\mathbf{k}}
\newcommand{\qv}{\mathbf{q}}
\newcommand{\pv}{\mathbf{p}}
\newcommand{\av}{\mathbf{a}}
\newcommand{\bv}{\mathbf{b}}

\title{High-symmetry points and paths in the Brillouin zone\\
       \large and what they dictate for the LiF polaron dispersion}
\author{Samuel Recker}
\date{\today}

\begin{document}
\maketitle

\begin{abstract}
The polaron-dispersion campaign measures $E_{0}(\pv)$ at nonzero external momentum and
extracts the effective mass $m^{*}$ from its curvature at $\Gamma$. Choosing \emph{which}
momenta to run means choosing directions in the Brillouin zone, and the conventional
choices --- ``$\Gamma\!\to\!X$'', ``$\Gamma\!\to\!L$'' --- come from crystal symmetry.
This note defines high-symmetry points and paths from scratch for the fcc lattice of
LiF, maps them onto the fractional-coordinate $n^{3}$ grid the sampler uses, derives what
cubic symmetry forces on the small-$\pv$ expansion of $E_{0}(\pv)$ (an isotropic scalar
$m^{*}$ for the non-degenerate conduction band, direction dependence only at quartic
order), and closes with an annotated reading list. The methodology note
(\texttt{latex/ab-initio-diagmc}) remains the authoritative reference for the sampler and
estimators; this note only covers the symmetry background.
\end{abstract}

\section{Reciprocal lattice, Brillouin zone, and fractional coordinates}
\label{sec:bz}

LiF crystallizes in the rock-salt structure: a face-centered-cubic (fcc) Bravais lattice
with a two-atom basis. With $a$ the conventional cubic lattice constant, the primitive
real-space vectors are
\begin{equation}
\av_{1}=\tfrac{a}{2}(0,1,1),\qquad
\av_{2}=\tfrac{a}{2}(1,0,1),\qquad
\av_{3}=\tfrac{a}{2}(1,1,0),
\end{equation}
and the reciprocal primitive vectors, fixed by $\bv_{i}\cdot\av_{j}=2\pi\delta_{ij}$, are
\begin{equation}
\bv_{1}=\tfrac{2\pi}{a}(-1,1,1),\qquad
\bv_{2}=\tfrac{2\pi}{a}(1,-1,1),\qquad
\bv_{3}=\tfrac{2\pi}{a}(1,1,-1).
\label{eq:recip}
\end{equation}
The first Brillouin zone (BZ) --- the Wigner--Seitz cell of the reciprocal lattice --- is
a truncated octahedron: eight regular hexagonal faces (perpendicular to the cube
diagonals $\langle111\rangle$) and six square faces (perpendicular to the cube axes
$\langle100\rangle$).

Throughout the codebase momenta are stored in \emph{fractional (crystal) coordinates}
$\mathbf{f}=(f_{1},f_{2},f_{3})$, meaning
\begin{equation}
\kv \;=\; f_{1}\bv_{1}+f_{2}\bv_{2}+f_{3}\bv_{3},
\label{eq:frac}
\end{equation}
on a $\Gamma$-centered $n^{3}$ grid $f_{i}=j_{i}/n$, $j_{i}\in\{0,\dots,n-1\}$. Two
consequences matter for everything below. First, the grid is closed under addition
modulo reciprocal-lattice vectors, which is what keeps every diagram's momentum
bookkeeping on-grid. Second, the map~\eqref{eq:frac} is \emph{not} an isometry --- the
$\bv_{i}$ are neither orthogonal nor unit-length --- so equal steps in fractional
coordinates are not equal steps in physical momentum, and any fit in $|\pv|$ must first
convert to Cartesian via~\eqref{eq:frac}.

\section{High-symmetry points}
\label{sec:points}

The point group of the crystal (for rock-salt: the full cubic group $O_{h}$, 48
operations --- rotations, mirrors, inversion) acts on reciprocal space and maps the BZ
onto itself. For a given $\kv$, the subgroup of operations that leave $\kv$ invariant
\emph{modulo a reciprocal-lattice vector} is called the \emph{little group} (or
stabilizer) of $\kv$. A generic $\kv$ has a trivial little group; points whose little
group is larger are called \emph{high-symmetry points} and carry conventional letters.
For the fcc BZ:

\begin{center}
\begin{tabular}{l l l l l}
\hline
Point & Fractional $(f_{1},f_{2},f_{3})$ & Cartesian $[\tfrac{2\pi}{a}]$ & Little group & Location \\
\hline
$\Gamma$ & $(0,0,0)$                     & $(0,0,0)$                     & $O_{h}$ (48)  & zone center \\
$X$      & $(\tfrac12,0,\tfrac12)$       & $(0,1,0)$                     & $D_{4h}$ (16) & square-face center, $\langle100\rangle$ \\
$L$      & $(\tfrac12,\tfrac12,\tfrac12)$& $(\tfrac12,\tfrac12,\tfrac12)$& $D_{3d}$ (12) & hexagon center, $\langle111\rangle$ \\
$W$      & $(\tfrac12,\tfrac14,\tfrac34)$& $(\tfrac12,1,0)$              & $D_{2d}$ (8)  & square--hexagon--hexagon corner \\
$K$      & $(\tfrac38,\tfrac38,\tfrac34)$& $(\tfrac34,\tfrac34,0)$       & $C_{2v}$ (4)  & hexagon--hexagon edge midpoint \\
$U$      & $(\tfrac58,\tfrac14,\tfrac58)$& $(\tfrac14,1,\tfrac14)$       & $C_{2v}$ (4)  & square--hexagon edge midpoint \\
\hline
\end{tabular}
\end{center}

\noindent
(Orders of the little groups in parentheses. The Cartesian columns follow
from~\eqref{eq:frac} with~\eqref{eq:recip}; e.g.\ for $X$:
$\tfrac12\bv_{1}+\tfrac12\bv_{3}=\tfrac{2\pi}{a}(0,1,0)$.) The set of all images of a
point under the point group is its \emph{star}; e.g.\ the $X$ star has the 6 members
$\tfrac{2\pi}{a}(\pm1,0,0),(0,\pm1,0),(0,0,\pm1)$, and the $W$ star (24 members) already
appears in the glossary because bands 6/7 are degenerate on it.

Why these points matter physically: any function with the lattice symmetry ---
$\varepsilon(\kv)$, $\omega_{\nu}(\qv)$, and the polaron dispersion $E_{0}(\pv)$ alike
--- satisfies $E(g\kv)=E(\kv)$ for every point-group operation $g$. At a high-symmetry
point the gradient components along symmetry-broken directions must vanish, so band
extrema, degeneracies, and critical points concentrate there. The conduction-band
minimum of LiF sits at $\Gamma$; the polaron ground state inherits that.

\section{Symmetry paths}
\label{sec:paths}

A \emph{symmetry path} is a piecewise-straight tour through high-symmetry points, along
edges and face normals of the irreducible BZ wedge. Band structures are functions on a
three-dimensional domain; the community convention is to plot them along a standard 1D
path so that any two calculations are visually comparable. For fcc the standard tour
(Setyawan--Curtarolo convention) is
\begin{equation*}
\Gamma \to X \to W \to K \to \Gamma \to L \to U \to W \to L \to K \,\big|\, U \to X,
\end{equation*}
where ``$|$'' marks a discontinuous jump. Nothing is physically special about the path
as a whole --- it is a plotting convention --- but the individual \emph{legs} from
$\Gamma$ are special: they are the extremal directions along which effective-mass
expansions are conventionally quoted, and symmetry constrains $E(\kv)$ most strongly
there. For the dispersion campaign only the first legs out of $\Gamma$ matter:
$\Gamma\to X$ (a $\langle100\rangle$ cube axis) and $\Gamma\to L$ (a $\langle111\rangle$
cube diagonal).

\section{The grid rays the sampler can run}
\label{sec:rays}

The external momentum \texttt{p\_ext} must be a point of the $n^{3}$ grid, so a
``path'' for us is a ray of grid points. The two legs from $\Gamma$:
\begin{align}
\Gamma\to X:\quad
\mathbf{f}_{j}=\tfrac{j}{n}(1,0,1)
\;\;\Rightarrow\;\;
\pv_{j}=\tfrac{2\pi}{a}\,\tfrac{2j}{n}\,(0,1,0),
&\qquad |\pv_{j}|=\tfrac{4\pi}{a}\,\tfrac{j}{n},
\qquad j=0,\dots,\tfrac{n}{2}, \\[2pt]
\Gamma\to L:\quad
\mathbf{f}_{j}=\tfrac{j}{n}(1,1,1)
\;\;\Rightarrow\;\;
\pv_{j}=\tfrac{2\pi}{a}\,\tfrac{j}{n}\,(1,1,1),
&\qquad |\pv_{j}|=\tfrac{2\pi\sqrt{3}}{a}\,\tfrac{j}{n},
\qquad j=0,\dots,\tfrac{n}{2}.
\end{align}
Both reach their endpoint exactly at $j=n/2$ (all ladder grids are even). Note the
change of appearance under the fractional$\to$Cartesian map: the ``diagonal-looking''
fractional ray $(j/n,0,j/n)$ is a Cartesian \emph{cube axis}, and the fractional body
diagonal is the Cartesian cube diagonal. The Cartesian step sizes differ between the
legs ($\tfrac{4\pi}{a}\tfrac1n$ vs $\tfrac{2\pi\sqrt3}{a}\tfrac1n$), which the fit
runbook must account for when comparing the two directions at matched $|\pv|$.

\section{What cubic symmetry dictates for $E_{0}(\pv)$ and $m^{*}$}
\label{sec:mass}

\paragraph{Assumptions, stated explicitly.} (i) The state being followed is
\emph{non-degenerate} at $\Gamma$ (true for the LiF conduction-band minimum, a single
Li-$s$-like band, and hence for the electron polaron built on it; the coupling
$g(\kv,\qv,\nu)$ respects the crystal symmetry, so the interacting dispersion inherits
$O_{h}$). (ii) $\pv$ stays below the phonon-emission threshold (below), so $E_{0}(\pv)$
is a sharp quasiparticle branch. (iii) Inversion symmetry (contained in $O_{h}$) plus
time reversal give $E_{0}(\pv)=E_{0}(-\pv)$, so the expansion is even.

Under these assumptions the most general small-$\pv$ expansion invariant under the cubic
group is
\begin{equation}
E_{0}(\pv)
= E_{0}
+ \frac{p_{x}^{2}+p_{y}^{2}+p_{z}^{2}}{2m^{*}}
+ A\,\bigl(p_{x}^{4}+p_{y}^{4}+p_{z}^{4}\bigr)
+ B\,\bigl(p_{x}^{2}+p_{y}^{2}+p_{z}^{2}\bigr)^{2}
+ \mathcal{O}(p^{6}).
\label{eq:expansion}
\end{equation}
The structure is forced by invariant theory: the only symmetric rank-2 tensor invariant
under $O_{h}$ is $\propto\delta_{ij}$ --- so the quadratic term is \emph{isotropic}, a
single scalar $m^{*}$ with no direction dependence --- while at quartic order there are
exactly two independent invariants, $\sum_{i}p_{i}^{4}$ and $|\pv|^{4}$. Restricted to
the two legs:
\begin{equation}
E_{0}(\pv)-E_{0}
= \frac{|\pv|^{2}}{2m^{*}} +
\begin{cases}
\bigl(A+B\bigr)\,|\pv|^{4} + \dots & \text{along } \Gamma\to X \;(\langle100\rangle),\\[4pt]
\bigl(\tfrac{A}{3}+B\bigr)\,|\pv|^{4} + \dots & \text{along } \Gamma\to L \;(\langle111\rangle),
\end{cases}
\label{eq:legs}
\end{equation}
using $\sum_{i}p_{i}^{4}=|\pv|^{4}$ on $\langle100\rangle$ and $=|\pv|^{4}/3$ on
$\langle111\rangle$. The two directions therefore agree on $m^{*}$ exactly and first
differ at quartic order, by $\tfrac{2}{3}A\,|\pv|^{4}$. Practical consequence: measuring
both legs is not new physics but a \emph{fit-window diagnostic} --- if quadratic fits
along $\Gamma\to X$ and $\Gamma\to L$ over the same $|\pv|$ window disagree beyond
combined errors, quartic contamination is leaking in and the window must shrink.

\paragraph{The degenerate contrast (why the hole is different).} The LiF valence-band
top at $\Gamma$ is threefold degenerate (F-$2p$-like). For a degenerate level the scalar
expansion~\eqref{eq:expansion} is replaced by a $3\times3$ matrix quadratic form (the
Luttinger--Kohn $\kv\!\cdot\!\pv$ Hamiltonian) with \emph{three} independent cubic
invariants (Luttinger parameters $\gamma_{1},\gamma_{2},\gamma_{3}$): heavy and light
branches, genuinely anisotropic and warped. This is precisely why the hole polaron
needs a multi-band, matrix-product-weight treatment and stays deferred (methodology
note, deferred-work section); nothing in this note's scalar analysis applies to it.

\paragraph{Where the quasiparticle branch ends.} A polaron at momentum $\pv$ can shed a
phonon and recoil; the dispersion flattens as $E_{0}(\pv)$ approaches the continuum edge
$E_{0}(\Gamma)+\omega_{\mathrm{LO}}$ ($\omega_{\mathrm{LO}}\approx0.08$~eV in LiF), and
beyond it there is no sharp branch to fit. All mass fits must stay well below this
threshold. Two sampling caveats tied to $|\pv|$: the projection gap that suppresses
excited-state contamination in the fixed-$\tau$ estimator shrinks as the threshold is
approached, so the required $\tau$ \emph{grows} with $|\pv|$ (hence a $\tau$-doubling
guard at the largest fitted momentum); and at $\pv\neq\Gamma$ the propagator phases are
genuinely complex, so $\langle\mathrm{sign}\rangle$ decays with $|\pv|$ and sets the
largest reachable grid index in practice.

\section{Reading}
\label{sec:reading}

\begin{description}
\sloppy
\item[Setyawan \& Curtarolo (2010).] W.~Setyawan, S.~Curtarolo, \emph{High-throughput
  electronic band structure calculations: Challenges and tools}, Comput.\ Mater.\ Sci.\
  \textbf{49}, 299 (2010); \href{https://arxiv.org/abs/1004.2974}{arXiv:1004.2974}.
  \emph{The} reference for BZ shapes, point labels, and standard paths for every Bravais
  lattice; the fcc figure and coordinate table cover everything in
  Section~\ref{sec:points}. Short and practical --- read this one first.
\item[Hinuma et al.\ (2017).] Y.~Hinuma, G.~Pizzi, Y.~Kumagai, F.~Oba, I.~Tanaka,
  \emph{Band structure diagram paths based on crystallography}, Comput.\ Mater.\ Sci.\
  \textbf{128}, 140 (2017); \href{https://arxiv.org/abs/1602.06402}{arXiv:1602.06402}.
  The modern, crystallographically careful successor (the SeeK-path paper). The
  companion web tool at
  \href{https://www.materialscloud.org/work/tools/seekpath}{materialscloud.org} draws
  the LiF BZ with $\Gamma$/$X$/$L$ labeled interactively --- the fastest way to build
  geometric intuition for Section~\ref{sec:points}.
\item[Yu \& Cardona.] P.~Y.~Yu, M.~Cardona, \emph{Fundamentals of Semiconductors},
  4th~ed., Springer (2010), ch.~2. The $\kv\!\cdot\!\pv$ method: why band extrema at
  high-symmetry points carry the mass-tensor structure their little group dictates,
  including the degenerate (Luttinger--Kohn) case of Section~\ref{sec:mass}.
\item[Dresselhaus, Dresselhaus \& Jorio.] \emph{Group Theory: Application to the
  Physics of Condensed Matter}, Springer (2008). The full machinery --- little groups,
  stars, representations at $\Gamma$/$X$/$L$, and $\kv\!\cdot\!\pv$ at high-symmetry
  points. The reference to reach for when the hole polaron (degenerate $\Gamma$ level)
  is picked up.
\item[Ashcroft \& Mermin.] \emph{Solid State Physics}, Saunders (1976), chs.~4--5
  and~8. Background on the reciprocal lattice, Brillouin zones, and Bloch bands
  (Section~\ref{sec:bz}), if a first-principles refresher is wanted.
\end{description}

\end{document}
